% Existing labels are retained to preserve all manuscript cross-references.
% The relocated single-robot routine is highlighted in full.
\begingroup\color{blue}
\section{\rev{Online Sensing and Plan Adaptation}}
\label{ch03:sec:partial}

\begin{revision}
The following single-robot routine applies the hard--soft synthesis of
Section~\ref{ch03:sec:soft-hard} after new sensing changes the workspace
abstraction. A revised continuation must remain compatible with the executed
history. Section~\ref{ch03:sec:multi-partial} reuses this routine after robots
exchange task-relevant region labels.
\end{revision}

In practice, the transition-system model may become inconsistent with the
actual workspace during execution. A plan synthesized off-line can therefore
become invalid, unsafe, or suboptimal after new information is acquired. This
motivates an online framework in which the robot monitors execution, updates the
workspace model, and revises the current plan while preserving validity and
safety. The transition system available at time $t \geq 0$ is
\begin{equation}\label{ch03:eq:t-FTS}
\mathcal{T}_c^t\triangleq (\Pi,\, \longrightarrow_c^t,\, \Pi_0,\,AP,\, L_c^t,\, w_c^t),
\end{equation}
where the region set $\Pi$, the initial-state set $\Pi_0$, and the
proposition set $AP$ are taken to be fixed. Only the transition relation, the
labeling function, and the transition weights are updated over time. The task specification is also fixed
during execution and retains the decomposition
$\varphi \triangleq
\varphi^{\textup{soft}} \wedge \varphi^{\textup{hard}}$
introduced in \eqref{ch03:eq:spec}. Let
$\tilde{\mathcal{A}}_{\varphi}$ denote the relaxed intersection automaton from
Definition~\ref{ch03:def:conj}.

\begin{problem}
Assume that the workspace is only partially known. The following questions
arise: (I) how should the robot's sensing information be modeled;
(II) how should the transition system be updated accordingly;
and (III) how can the motion and task plan be kept valid and safe throughout execution?
\end{problem}

\subsection{\rev{Initial plan synthesis}}

At $t=0$, the initial plan is synthesized by
Section~\ref{ch03:sec:soft-hard} from $\mathcal{T}_c^0$ and~$\varphi$. For
$t \geq 0$, let $R^t$ be the current safe accepting run,
$\tau^t \triangleq R^t|_{\Pi}$ the corresponding plan, and
$\tilde{\mathcal{A}}_r^t$ the associated product automaton, generated on the
fly. Even if the soft component is initially infeasible,
Theorem~\ref{ch03:theo:safety} guarantees that $\tau^0$ is valid and safe with
respect to $\mathcal{T}_c^0$.

\subsection{\rev{Real-time knowledge update}}
\label{ch03:subsec:knowledge-update}

The robot acquires workspace information through sensing during execution.
Only locally observed information is used here. Inter-robot information
exchange is introduced in Section~\ref{ch03:sec:multi-partial}.
\begin{revision}
Let $\textbf{Sense}^t$ denote the sensing information obtained
at time $t \geq 0$, either upon reaching a region or while traversing a
transition:
\begin{equation}
\textbf{Sense}^t \triangleq
\{\, ((\pi,\, S,\, S_{\neg}),\, E,\, E_{\neg}) \,\},
\label{ch03:eq:sense-info}
\end{equation}
where $\pi \in \Pi$ is an observed region, $S$ and $S_{\neg}$ contain
propositions confirmed true and false at~$\pi$, respectively,
$(\pi_i,\, \pi_j,\, w) \in E$ proposes a transition or a weight update, and
$(\pi_i,\, \pi_j) \in E_{\neg}$ identifies a transition to remove.
\end{revision}
Thus $\textbf{Sense}^t$ encodes
the symbolic workspace information revealed at time~$t$. The extraction of
\eqref{ch03:eq:sense-info} from raw sensor data depends on the sensing modality
and is not needed here.

\subsection{\rev{Transition system update}}\label{ch03:subsubsec:updateT}

At time~$t$, the robot may receive new information through
$\textbf{Sense}^t$. This information is used to update the current
transition-system model. Let $\mathcal{T}_c^{t^-}$ and
$\mathcal{T}_c^{t^+}$ denote the transition systems immediately before and
after the update at time~$t$, respectively.
If $(\pi,\, S,\, S_{\neg})$ is obtained from \eqref{ch03:eq:sense-info}, then
the labeling function is updated according to
$L_c^{t^+}(\pi) \triangleq
\big(L_c^{t^-}(\pi) \cup S\big) \setminus S_{\neg}$.
\begin{revision}
An online update adds an edge only after identifying a local controller
that satisfies Definition~\ref{ch02:def:transition} under the updated
workspace information. The controller is stored with the edge. An
existing edge can be reweighted while its controller remains feasible.
The update removes the edges in $E_{\neg}$ and signals an interruption
if it removes the active edge or invalidates its controller.
\end{revision}
For convenience, the set of regions whose outgoing adjacency information
must be reconstructed is defined by
\begin{equation}
\widehat{\Pi} \triangleq
\widetilde{\Pi} \cup
\{\, \pi_i \mid (\pi_i,\, \pi_j,\, w) \in E
\text{ or } (\pi_i,\, \pi_j) \in E_{\neg} \,\},
\label{ch03:eq:widehat}
\end{equation}
where $\widetilde{\Pi} \subseteq \Pi$ denotes the set of regions whose labeling
function changes during the update. If $\textbf{Sense}^t$ is empty, then
$\mathcal{T}_c^{t^+} = \mathcal{T}_c^{t^-}$.

\subsection{\rev{Safety-ensured plan revision}}\label{ch03:subsec:realtime}

Because $\mathcal{T}_c^t$ may change during execution, the current plan must
be re-evaluated for validity, safety, and optimality. The update first modifies
the product automaton representation and then decides whether the current run
can be retained, locally repaired, or replaced.

(I) \emph{Product automaton update}.
Let $\tilde{\mathcal{A}}_r^{t^-}$ and $\tilde{\mathcal{A}}_r^{t^+}$ denote the
product automata before and after the update. Reconstructing
$\tilde{\mathcal{A}}_r^{t^+}$ globally, or re-evaluating every affected
transition as in~\citep{guo2013revising}, is often unnecessary. Instead, update
information is incorporated lazily into the adjacency routine: product states
whose workspace component lies in $\widehat{\Pi}$ have their outgoing
transitions rebuilt only when revisited. The remaining questions are whether
$\tau^t$ is still valid and safe, and how it should be revised if not.

(II) \emph{Verifying validity and safety}.
The prefix-suffix representation of $R^t$ contains finitely many distinct
transitions. Let $\texttt{Edge}(R^t)$ be the set of transitions appearing in
that representation. Scanning this set determines whether the current plan
remains valid and safe.

\begin{definition}\label{ch03:def:vsTS}
Given the updated transition system $\mathcal{T}_c^{t^+}$, a transition
$(\langle \pi_i,\, q_m \rangle,\,
\langle \pi_j,\, q_n \rangle) \in \textup{\texttt{Edge}}(R^t)$ is called
\emph{invalid} if $(\pi_i,\, \pi_j) \notin \longrightarrow_c^{t^+}$, and
\emph{unsafe} if
$L_c^{t^+}(\pi_i) \notin \chi_1(q_m|_{Q_1},\, q_n|_{Q_1})$.
\end{definition}

Let $\Xi$ and $\aleph$ denote the sets of invalid and unsafe transitions of
$R^t$, respectively. These sets are obtained by scanning the transitions in
$\texttt{Edge}(R^t)$ and checking whether a transition has been removed from
$\mathcal{T}_c^{t^+}$ or whether an updated label violates the hard-automaton
constraint.

\begin{theorem}\label{ch03:theo:validity}
Assume that $R^t$ is an accepting run of $\tilde{\mathcal{A}}_r^{t^-}$ and that
$\mathcal{T}_c^{t^-}$ is updated to $\mathcal{T}_c^{t^+}$. Let $\Xi$ and
$\aleph$ be the invalid and unsafe transitions above. Then the
following statements hold:
(I) $\tau^t$ remains valid if and only if $\Xi = \emptyset$;
(II) $\tau^t$ remains safe if $\Xi = \emptyset$ and $\aleph = \emptyset$.
\end{theorem}

\begin{proof}
Since $R^t$ is an accepting run of $\tilde{\mathcal{A}}_r^{t^-}$,
Theorem~\ref{ch03:theo:safety} implies that $\tau^t$ is valid and safe for
$\mathcal{T}_c^{t^-}$. Statement~(i) follows directly from the definition of
$\Xi$: the plan remains valid precisely when none of its transitions has been
removed from the updated transition relation.
For statement~(ii), assume that $\Xi = \emptyset$ and $\aleph = \emptyset$.
Then every transition of $R^t$ remains valid in $\mathcal{T}_c^{t^+}$ and
still satisfies the hard-automaton constraint under the updated labeling
function. Hence $R^t$ remains an accepting run of
$\tilde{\mathcal{A}}_r^{t^+}$, and Theorem~\ref{ch03:theo:safety} yields that
$\tau^t$ remains safe.
\end{proof}

(III) \emph{Plan revision}:
If either $\Xi$ or $\aleph$ is nonempty, the current plan may no longer be
executable or safe. Replanning from the current region alone can yield a valid
future plan, but it may be incompatible with the product-state history induced
by the executed trajectory. Let $\tau_{\textup{past}}$ and
$R_{\textup{past}}$ denote the executed trajectory and its recorded product run,
and let
$\tau_{\textup{comp}} \triangleq
\tau_{\textup{past}} + \tau_{\textup{new}}$. Since the B\"uchi automaton is
nondeterministic, several product states may correspond to the same current
workspace region, which leads to the following compatibility problem.

\begin{problem}\label{ch03:prob:concatenate}
Assume that $\Xi$ or $\aleph$ is nonempty. Given the executed past trajectory
$\tau_{\textup{past}}$, find a new plan $\tau_{\textup{new}}$ such that the
complete trajectory
$\tau_{\textup{comp}} = \tau_{\textup{past}} + \tau_{\textup{new}}$
is valid and safe.
\end{problem}

To formalize the compatibility constraint, let
$\bar{\tau}_{\textup{past}} \triangleq
\tau_{\textup{past}}[1:(|\tau_{\textup{past}}|-1)]$, that is, the executed
trajectory without its last state. Because $\tilde{\mathcal{A}}_{\varphi}$ is
nondeterministic, the trace of $\bar{\tau}_{\textup{past}}$ may induce several
finite runs in $\tilde{\mathcal{A}}_{\varphi}$. Denote this set by
$r_{\tau_{\textup{past}}} \triangleq
\{\, r_{\bar{\sigma}}
\mid \bar{\sigma} = \texttt{trace}(\bar{\tau}_{\textup{past}}) \,\}$.
The corresponding set of finite runs in the updated product
automaton is
$R_{\tau_{\textup{past}}} \triangleq
\{\, R \mid R|_{\Pi} = \tau_{\textup{past}},\;
R|_{Q} \in r_{\tau_{\textup{past}}} \,\}$.
Accordingly, the product states consistent with the executed past is
\begin{equation}
Q'_{\tau_{\textup{past}}} \triangleq
\{\, \texttt{last}(R) \mid R \in R_{\tau_{\textup{past}}} \,\},
\label{ch03:eq:set-past}
\end{equation}
where $\texttt{last}(R)$ is the last state of the finite run~$R$. By
construction, $R_{\textup{past}} \in R_{\tau_{\textup{past}}}$ and the recorded
current product state belongs to $Q'_{\tau_{\textup{past}}}$. The revision
algorithm first attempts to repair the current run by inserting a bridging
segment around each invalid or unsafe transition, following the spirit of
\citet{guo2013revising}. \rev{If no bridge exists, the revision calls
\texttt{OptRun} with start set $Q'_{\tau_{\textup{past}}}$ and the same suffix
weight $\gamma$ used in the original synthesis.}
\rev{To apply the penalty bound in \eqref{ch03:underalpha} to exact
resynthesis after an update, compute the feasible continuation in the
updated unrelaxed hard--soft product from a state reachable along the
executed history in that product. Evaluate its motion cost with the
same $\gamma$ used in the relaxed continuation search.}

\begin{algorithm}[t]
\caption{Revise the current plan, \texttt{Revise}()}
\label{ch03:alg:revise}
\DontPrintSemicolon
\KwIn{\rev{$\Xi$, $\aleph$, $R^{t^-}$, $\tilde{\mathcal{A}}_{r}^{t^+}$,
      $Q'_{\tau_{\textup{past}}}$, $\gamma>0$}}
\KwOut{$R^{t^+}$}

\tcc{\rev{Try a local bridge; fall back to synthesis consistent with the executed past.}}
\ForAll{$(q'_s,\, q'_{s+1}) \in (\Xi \cup \aleph)$}{
  \If{$(q'_s,\, q'_{s+1}) \in R^{t^-}_{\textup{pre}}$}{
    $bridge = \texttt{DijksTarget}(
    \tilde{\mathcal{A}}_{r}^{t^+},\, q'_{s-1},\,
    \texttt{tail}(q'_s,\, R^{t^-}_{\textup{pre}}))$\;

    \If{$bridge \neq \emptyset$}{
      $R^{t^-}_{\textup{pre}} =
      \texttt{head}(q'_{s-1},\, R^{t^-}_{\textup{pre}})
      + bridge
      + \texttt{tail}(\texttt{last}(bridge),\, R^{t^-}_{\textup{pre}})$\;
    }
    \Else{
      \rev{$R^{t^+} =
      \texttt{OptRun}(\tilde{\mathcal{A}}_{r}^{t^+},
      Q'_{\tau_{\textup{past}}},\gamma)$}\;
      \textbf{return} $R^{t^+}$\;
    }
  }

  \If{$(q'_s,\, q'_{s+1}) \in R^{t^-}_{\textup{suf}}$}{
    Apply the same procedure to $R^{t^-}_{\textup{suf}}$\;
  }
}
\textbf{return}
$R^{t^+} = \langle R^{t^-}_{\textup{pre}},\, R^{t^-}_{\textup{suf}} \rangle$\;
\end{algorithm}

\begin{theorem}\label{ch03:theo:new-safe}
Assume that there exists an accepting run of
$\tilde{\mathcal{A}}_{r}^{t^+}$ whose projection onto $\Pi$ extends the
executed past trajectory $\tau_{\textup{past}}$. Then
Algorithm~\ref{ch03:alg:revise} returns a continuation
$\tau_{\textup{new}}$ such that
$\tau_{\textup{comp}} = \tau_{\textup{past}} + \tau_{\textup{new}}$
is valid and safe.
\end{theorem}

\begin{proof}
If a compatible accepting continuation exists, then there is an accepting run
of $\tilde{\mathcal{A}}_{r}^{t^+}$ whose initial state belongs to
$Q'_{\tau_{\textup{past}}}$. Algorithm~\ref{ch03:alg:revise} first attempts to
repair the current run locally by inserting bridging segments around the
transitions in $\Xi \cup \aleph$. Whenever such a bridge is found, the modified
run remains consistent with the unchanged portion of the original accepting
run.
\rev{If no bridge can be found for some transition, the algorithm calls
\texttt{OptRun} on the updated product with start set
$Q'_{\tau_{\textup{past}}}$ and suffix weight $\gamma$.}
By assumption, a compatible accepting run exists,
so this search succeeds and returns one such run. In either case, the resulting
run is an accepting run of $\tilde{\mathcal{A}}_{r}^{t^+}$ whose projection
extends $\tau_{\textup{past}}$. Theorem~\ref{ch03:theo:safety} then implies
that the complete trajectory $\tau_{\textup{comp}}$ is valid and safe.
\end{proof}

Algorithm~\ref{ch03:alg:scheme} later combines the revision procedure of
Algorithm~\ref{ch03:alg:revise} with the optimal search routine of
Algorithm~\ref{ch02:alg:optimal}; additional implementation details appear
in~\citep{guo2014cooperative}. \rev{At a confirmed admissible state, the
supervisor uses the revised run to select the next target. If an update
invalidates the active controller, the supervisor handles the interruption
before further symbolic execution.}

%%%%%%%%%%%%%%%%%%%%%%%%%%%%%%%%%%%%%%%%%%%%%
%%%%%%%%%%%%%%%%%%%%%%%%%%%%%%%%%%%%%%%%%%%%%
\endgroup
